% !TEX program = lualatex
\documentclass[11pt,a4paper]{article}

\usepackage{fontspec}
\usepackage[{{babel}}]{babel}
\usepackage{unicode-math}
% {{fr:Libertinus : belles polices libres avec lettres grecques|en:Libertinus: nice free fonts with Greek letters}}
\setmainfont{Libertinus Serif}
\setsansfont{Libertinus Sans}
\setmathfont{Libertinus Math}
\usepackage{microtype}
\usepackage[margin=2.5cm]{geometry}
\usepackage{amsmath}
\usepackage{graphicx}
\usepackage{booktabs}
\usepackage{siunitx}
\usepackage{csquotes}
\usepackage[backend=biber, style=authoryear]{biblatex}
\addbibresource{references.bib}
\usepackage[hidelinks]{hyperref}
\usepackage[{{babel}}, nameinlink]{cleveref}

\title{{{title}}}
\author{{{author}}}
\date{{{date}}}

\begin{document}
\maketitle

\section{Introduction}
\label{sec:intro}

{{fr:Avec LuaLaTeX, on écrit directement en Unicode : α, β, γ et même « guillemets ».|en:With LuaLaTeX you write Unicode directly: α, β, γ and even “quotes”.}}
{{fr:Les unités sont gérées par siunitx : \qty{9.81}{\metre\per\second\squared}.|en:Units are handled by siunitx: \qty{9.81}{\metre\per\second\squared}.}}
{{fr:Knuth a créé \TeX{} \autocite{knuth1984}.|en:Knuth created \TeX{} \autocite{knuth1984}.}}

\begin{equation}
  ∇ ⋅ 𝐄 = \frac{ρ}{ε_0}
  \label{eq:gauss}
\end{equation}

{{fr:\Cref{eq:gauss} est l'une des équations de Maxwell (voir \cref{sec:intro}).|en:\Cref{eq:gauss} is one of Maxwell's equations (see \cref{sec:intro}).}}

\printbibliography
\end{document}
